\documentclass[11pt]{article}
\usepackage[utf8]{inputenc}\usepackage[T1]{fontenc}\usepackage[margin=1in]{geometry}
\usepackage{mathptmx}\usepackage{booktabs}\usepackage{tabularx}\usepackage{enumitem}\usepackage{microtype}
\usepackage[colorlinks=true,linkcolor=blue!60!black,urlcolor=blue!60!black]{hyperref}\usepackage{fancyhdr}
\pagestyle{fancy}\fancyhf{}\fancyhead[L]{\small AI-Level --- The Free Edition}\fancyhead[R]{\small Yang}\fancyfoot[C]{\thepage}
\renewcommand{\headrulewidth}{0.3pt}\setlength{\headheight}{14pt}\setlength{\parskip}{4pt}
\begin{document}
\begin{center}
{\LARGE\bfseries AI-Level}\\[6pt]
\end{center}

\section*{From LLMs to a Type II Civilization}

\subsection*{The Five Circles and the Grand Problems of Our Time}

\textbf{Chengshuai Yang}

\textbf{The Free Edition · Version 3 · September 2026}

\emph{Short. Plain. Free to read and share. The full book holds the theory, the measurements and a card for every problem named here.}

\medskip\noindent\rule{\textwidth}{0.3pt}\medskip

\section*{Start with a pump}

Ask an AI to design a machine that repairs a village water pump. It gives you a good answer in ten seconds. Now what?

Someone has to make the parts. Someone has to assemble them, test the machine on a real broken pump, and keep it working next year. If all of that happens, a community gains something it did not have before: it can fix its own water. If any link fails, the ten-second answer was worth nothing.

That chain --- from a good answer to a capability people can depend on --- is what this book is about. Every link has a question hiding in it: what was actually built, who checked it, and what still stands in the way. The most important link is the one people skip: \textbf{who checked it, and could the thing that made it have checked itself?}

\section*{What "smart" hides}

Everyone asks how smart AI is. It is the wrong question, because it hides seven different ones:

\begin{itemize}[leftmargin=1.4em,itemsep=2pt]
\item \textbf{Can it solve the problem in front of it?} (cognition)
\item \textbf{Does it remember, and does remembering change what it does next?} (memory and learning)
\item \textbf{Can it improve its own methods, and prove the improvement?} (individual growth)
\item \textbf{Can it work inside a group where someone else checks the work?} (organization)
\item \textbf{How hard a job can you hand it?} (task difficulty)
\item \textbf{How much thinking does a person still have to supply?} (human help)
\item \textbf{Does it know what it knows, and say when it can't?} (self-awareness --- knowing its limits, not consciousness)
\end{itemize}

Two rules make the answers honest.

\textbf{A confident wrong answer costs more than "I can't."} Most AI tests give a point for right and nothing for wrong, so guessing is free. We charge for it. An AI that hands you wrong work as if it were finished is worse than one that tells you it is stuck --- and a test that punishes bluffing changes which AI looks best.

\textbf{The AI is measured together with its scaffolding.} The same model behaves differently depending on the tools, memory and checks wrapped around it. A score belongs to the pair. We say which pair.

Measured this way, today's language models are at the entrance. They reason well. Under an ordinary API call they remember nothing once the conversation ends. And in every system anyone has measured, the two decisions that matter --- is this work right, and shall we use it --- are still made by a person. That is not an insult. It is where the road begins.

\section*{The loop, and the five circles}

Every real improvement, in any field, has the same four steps:

\begin{enumerate}[leftmargin=1.6em,itemsep=2pt]
\item \textbf{Identify} what to change.
\item \textbf{Implement} the change.
\item \textbf{Validate} --- check that it worked, honestly.
\item \textbf{Deploy} --- put it into use.
\end{enumerate}

A circle is defined by \emph{what kind of thing} that loop can change.

\begin{center}\small\begin{tabularx}{\textwidth}{@{}p{2.6cm}XX@{}}	oprule Circle & What the loop changes & How fast you find out if it worked \\ \midrule
\textbf{I --- Software} & programs, including the AI's own & seconds to hours \\
\textbf{II --- Hardware} & the chips it runs on & design: days; manufacture: months \\
\textbf{III --- Machines} & robots and factories that build and repair themselves & months to years \\
\textbf{IV --- Life} & biology --- treatments for cancer, aging, damaged organs, the brain & years; bodies have their own clock \\
\textbf{V+ --- Space} & industry in orbit that grows by itself & decades \\
\bottomrule\end{tabularx}\end{center}

The last circle ends at a \textbf{Dyson swarm}: thousands of collectors orbiting the Sun, built by an industry that grows from the energy it captures. A civilization that uses energy on that scale is what astronomers call \textbf{Type II}. (Careful: "Circle II" is chips; "Type II" is a star's worth of energy. Different things.)

Two facts shape everything else.

\textbf{The thing being improved never checks itself.} A student cannot grade their own exam. A system that approves its own changes has proved nothing. So in every circle, two of the four steps need an outside party --- a referee. Building a referee that cannot be fooled is the first grand problem, and every circle inherits it.

\textbf{Smarter alone does not get you out.} Designing a chip, a machine or a solar collector is ordinary engineering thinking. What is hard about the outer circles is everything \emph{after} the design: making it, testing it in the world, keeping it running for years. Those are problems of materials, time and organization. The outer circles are won by organizations, not by one brilliant mind --- and the last circle, sixteen light-minutes across, cannot be run by any single mind at all.

The circles overlap and develop in parallel. Biology does not wait for robots. Chips and factories improve each other. What changes is where the \emph{bottleneck} sits, and that is what to work on.

\section*{Forty problems, circle by circle}

Under several of the forty sit eighteen more specific research problems --- cancer as an evolving population, a room-temperature superconductor, the shape of spacetime from entanglement. Each has a precise statement of what "solved" would mean, and none has a benchmark yet. That is the honest state of every problem in this book, and an invitation.

\subsection*{Circle I --- Software}

\begin{itemize}[leftmargin=1.4em,itemsep=2pt]
\item \textbf{I-01 Memory that learns.} Keep experience across restarts and get better at \emph{new} tasks because of it --- not just store text.
\item \textbf{I-02 Dependable long work.} Finish long tasks, recover from errors, notice when something is missing.
\item \textbf{I-03 Self-improvement, proven.} Change its own methods in a way an outside check confirms helped, without breaking what already worked.
\item \textbf{I-04 Better at getting better.} Improve the process that finds and tests improvements.
\item \textbf{I-05 Discovery that becomes ability.} Find something new, confirm it, and be measurably better afterwards.
\item \textbf{I-06 A referee that cannot be fooled.} Tests the AI never sees, checks it cannot rewrite, and the rule that charges for bluffing.
\item \textbf{I-07 Organizations of AIs.} Real division of labour, where the one proposing a change is never the one approving it.
\item \textbf{I-08 A ruler that does not move.} A way to score AI progress that means the same thing next year.
\end{itemize}

\textbf{Four physics problems an AI can attack with a computer and public data:} the shape of spacetime from quantum entanglement (\emph{quantum gravity}); what makes galaxies spin the way they do (\emph{dark matter}); whether the force stretching the universe is changing (\emph{dark energy}); whether black holes destroy information. A good result on any of these is a result \emph{inside a model} --- the universe still gets the last word --- but they are the cleanest tests of whether an AI can discover something and be checked by data it cannot see.

\subsection*{Circle II --- Hardware}

\begin{itemize}[leftmargin=1.4em,itemsep=2pt]
\item \textbf{II-01 Chip designs that work when built.}
\item \textbf{II-02 Factories that run themselves} through the steps that turn a design into a device.
\item \textbf{II-03 The yield loop:} find defects, correct drift, raise the share of chips that work --- at the speed of physics.
\item \textbf{II-04 Better devices and connections.}
\item \textbf{II-05 More computation per watt.}
\item \textbf{II-06 Cooling and staying reliable} for years.
\item \textbf{II-07 Keeping the tools and materials coming.}
\item \textbf{II-08 Running on what you made.} The new hardware goes into the system's own computers.
\end{itemize}

Five more specific problems here: quantum computers that stay correct as they grow; the control pulses that run them; compilers that fit a program onto real quantum hardware; single molecules as electronic parts; a room-temperature superconductor. Each is a \emph{route}, not a requirement; a prediction is not a device, and a device is not yet a computer.

\subsection*{Circle III --- Machines and matter}

\begin{itemize}[leftmargin=1.4em,itemsep=2pt]
\item \textbf{III-01 Robots that handle the unplanned.}
\item \textbf{III-02 Diagnose, repair, verify.}
\item \textbf{III-03 From rock to product:} mining, refining and manufacturing as one chain.
\item \textbf{III-04 Precision} --- the parts, instruments and electronics complex machines need.
\item \textbf{III-05 Supplies that keep flowing} through disruption.
\item \textbf{III-06 A machine that builds a working copy of itself} --- and lists honestly every part it had to import.
\item \textbf{III-07 Energy and the services production runs on.}
\item \textbf{III-08 Checking with instruments the machine does not control.}
\end{itemize}

Two more specific problems here: building matter one atom at a time; DNA that folds itself into designed shapes. Both are precision manufacture at the smallest scale. Neither is yet a factory that makes factories.

\subsection*{Circle IV --- Life}

\begin{itemize}[leftmargin=1.4em,itemsep=2pt]
\item \textbf{IV-01 Cause, not correlation:} which mechanisms drive an outcome, tested by changing them.
\item \textbf{IV-02 Labs that learn} --- automated experiments whose results change what the lab tries next, and get repeated by someone else.
\item \textbf{IV-03 Cancer.} Treat a tumour as an evolving population and keep it controllable instead of breeding resistance.
\item \textbf{IV-04 Aging.} Turn back a cell's biological age without losing its identity; slow the rate of aging itself.
\item \textbf{IV-05 Regeneration.} Steer a wound toward repair instead of scar.
\item \textbf{IV-06 The brain.} Predict and reverse the protein clumping behind Alzheimer's and Parkinson's.
\item \textbf{IV-07 Delivery.} Drugs that go only where they are needed.
\item \textbf{IV-08 Benefit that reaches people} and keeps being measured after it does.
\end{itemize}

This circle has a clock nobody can hurry. A treatment is tested in real bodies over real time, with consent and accountability. A younger reading on an aging clock is a clue, not a cure; less protein clumping is a mechanism, not recovery. The discovery loop can run fast inside a validation loop that cannot --- and that is fine, because health is a purpose in itself, not a step toward space.

\subsection*{Circle V+ --- Space}

\begin{itemize}[leftmargin=1.4em,itemsep=2pt]
\item \textbf{V-01 Supplies from off Earth.}
\item \textbf{V-02 Industry in orbit that grows faster than it wears out.} Not one collector: a loop that builds collectors from the energy it collects.
\item \textbf{V-03 Collectors that last.}
\item \textbf{V-04 Useful energy} --- converting, storing, moving it.
\item \textbf{V-05 Getting rid of the heat.} Every watt used must be radiated away.
\item \textbf{V-06 Orbits and collisions.}
\item \textbf{V-07 Repair over decades}, with help far away.
\item \textbf{V-08 Organization across light-minutes.} No single mind, human or machine, can run a structure the size of Earth's orbit in real time. This is the deepest problem in the book, and it is a problem of organization, not intelligence.
\end{itemize}

Quantum networks may one day carry trust across distance (they do not carry messages faster than light). But the eight problems above are the ones that decide this circle, and they are the least worked-on problems in this book. That gap is worth filling.

\section*{Where to begin}

Not "which circle is humanity in." Which bottleneck is in front of \emph{you}.

\begin{itemize}[leftmargin=1.4em,itemsep=2pt]
\item \textbf{Your AI already proposes useful changes?} Work on the checking --- the referee, I-06. That is where most progress is being lost right now.
\item \textbf{It doesn't?} Work on the reasoning, the memory, or the ability to finish long tasks first.
\item \textbf{You run a factory or a fab?} Find the one component, process or repair skill that stops the next machine from being built. Fix that; then see what stops the one after.
\item \textbf{You run a lab?} Find the weakest link between your hypothesis and a person getting better --- usually a measurement you cannot yet trust.
\item \textbf{You think about space?} Supplies and construction first; heat and repair next; coordination when delay actually starts to bind.
\end{itemize}

For a small team with a computer, the shortest loops in this book are the four physics problems and the four quantum-software problems: computation and calibration, no clinic and no fab. One of them could move up a rung this year.

\section*{Which problem is yours}

\begin{itemize}[leftmargin=1.4em,itemsep=2pt]
\item \textbf{A student:} reproduce one published result on one of the physics or quantum problems, then build the test that would let an AI be checked against it. That test does not exist yet.
\item \textbf{A software engineer:} run a memory experiment with the memory switched off as the control. An afternoon of work; almost nobody publishes it.
\item \textbf{A laboratory:} pick one assay and make it reproducible enough that another lab gets the same number.
\item \textbf{A manufacturer:} publish the full list of what your most automated line still imports. The list is the map.
\item \textbf{A funder:} pay for referees and replications, not only for builders and discoveries.
\item \textbf{A community:} name the pump. The problem that matters to you is a grand problem if solving it creates a capability others can use.
\end{itemize}

\section*{What "solved" looks like}

Before you start, write five lines:

\begin{enumerate}[leftmargin=1.6em,itemsep=2pt]
\item \textbf{Purpose.} Whose situation improves if this works?
\item \textbf{Constraint.} What is actually stopping it --- and what is the evidence?
\item \textbf{Contribution.} The specific thing you will build.
\item \textbf{Test.} What counts as progress, what counts as failure, and who --- not you --- will check.
\item \textbf{Next use.} How the next person picks it up.
\end{enumerate}

Then build the smallest thing that would pass the test. Get it checked. Hand it on.

\textbf{Choose a problem that matters. Build a capability that helps solve it. Show that it works to someone who is not you. Make it usable by the next person.}

\section*{Where this goes next}

\begin{itemize}[leftmargin=1.4em,itemsep=2pt]
\item \textbf{The full book:} \emph{AI-Level: From LLMs to a Type II Civilization} --- the theory of the seven measurements and the five circles; the instrument that measures an AI and its scaffolding together, what it has measured, and what it gets wrong; each circle in depth with the evidence for its bottleneck and the test that would prove us wrong; a full card for every one of the forty problems; and what happens to human work as each circle closes.
\item \textbf{The sources:} the circle framework and the AI-Level papers behind this book are public, and the full book cites them.
\end{itemize}

\emph{This free edition may be copied and shared without permission. The full book supports the research it describes.}

\end{document}
